\par\noindent\textbf{(1)} 单态投影P\_s的本征值为1及三重零，混合单位阵后，ρ的本征值为 (1+3p)/4 和三重 (1−p)/4。对0≤p≤1均非负、迹为1；极限p=0为最大混态、p=1为纯singlet。
\par\noindent\textbf{(2)} singlet投影在第二比特部分转置后本征值为三重1/2和单个−1/2。于是ρ\textasciicircum{}\{T\_B\}的本征值为三重(1+p)/4和单个(1−3p)/4。两qubit PPT判据为充分必要条件，因此p>1/3出现负部分转置、态纠缠，p≤1/3可分。
\par\noindent\textbf{(3)} Werner态经spin-flip的Wootters并合构造，concurrence C=max(0,(3p−1)/2)。纯度Trρ²=[(1+3p)²+3(1−p)²]/16=(1+3p²)/4。它们分别捕捉纠缠强度和相干混合程度，二者阈值不能混同。
\par\noindent\textbf{(4)} 相关矩阵 T\_\{ij\}=Trρ(σ\_i⊗σ\_j)=−pδ\_ij，故Horodecki准则 CHSH最大值 S\_max=2√(m₁+m₂)=2√2p。仅p>1/√2违反局域CHSH界2，因此存在1/3<p≤1/√2的纠缠却未违反该Bell不等式的区域。
